% GENERATED FILE. Edit canonical lessons, then run npm run generate:books.
% canonical-source-hash: fnv1a64:8a387cfc20b385d5
% canonical-lessons: ZH-C443-qingche, ZH-C443-baohu, ZH-W443-hu-guard, ZH-C443-baohuqu, ZH-C443-husong, ZH-C443-buyu

\chapter{Keeping It Clean}
\label{ch:zh-keeping-it-clean}
\hlchaptermodality{443}

\begin{chapteropening}
\textbf{By the end of this chapter:} \emph{I can say the water is clear, protect a river, visit a nature reserve, escort someone and talk about fishing.}

Complete the last lesson of chapter 443: I can say the water is clear, protect a river, visit a nature reserve, escort someone and talk about fishing.
\end{chapteropening}

\section[qīngchè]{qīngchè — crystal-clear (of water)}

\label{lesson:ZH-C443-qingche}



\begin{quote}
Before the new one: say the Mandarin for to rise after rain, then the Mandarin for a great river.
\end{quote}

\subsection*{You'll want to know: qīngchè}
\emph{qīngchè} — \textquotedblleft{}crystal-clear (of water)\textquotedblright{}.

\begin{quote}
\emph{zhè tiáo hé de shuǐ hěn qīngchè} — The water in this river is crystal clear.
\end{quote}

Its characters come later; for now it is a word for the ear.

\begin{cousinweb}
\emph{qīngchè} is built from \emph{qīng} 'clear' and \emph{chè} 'clear water', both inherited from Middle Chinese, \emph{tshjeng} and \emph{drjet} (Guangyun, in Baxter's transcription), which glosses the second as water that has settled clear.
\end{cousinweb}

\begin{grammarlens}[title={What you've built}]
One more word for getting around.
\end{grammarlens}

\subsection*{Your turn}
\begin{itemize}
\raggedright
  \item \emph{Say it:} \emph{qīngchè}
  \item \emph{Say it:} \emph{qīngchè}, once more
  \item \emph{Say it:} \emph{zhè tiáo hé de shuǐ hěn qīngchè} — The water in this river is crystal clear.
  \item \emph{Recall:} say the Mandarin for to rise after rain, then the Mandarin for a great river, then say \emph{qīngchè} again
\end{itemize}

\begin{tcolorbox}[breakable,colback=teal!4,colframe=teal!35!black,title={Before you move on}]
What does qīngchè mean? (\textquotedblleft{}Crystal-clear (of water)\textquotedblright{}.) Say it once more.
\end{tcolorbox}
\section[\texorpdfstring{\zh{保护} (bǎohù)}{保护 (bǎohù)}]{\zh{保护} (bǎohù) — to protect; to keep from harm}

\label{lesson:ZH-C443-baohu}



\begin{quote}
Before the new one: say the Mandarin for a great river, then the Mandarin for crystal-clear.
\end{quote}

\subsection*{You'll want to know: \zh{保护}}
\textbf{\zh{保护}} — \emph{bǎohù} — \textquotedblleft{}to protect; to keep from harm\textquotedblright{}.

\begin{quote}
\textbf{\zh{我们应该保护河水}} \emph{wǒmen yīnggāi bǎohù héshuǐ} — We should keep the river water clean.
\end{quote}

The new character, in pieces: \zh{护} — seven strokes in two pieces: on the left the hand piece of \zh{打}, three strokes; on the right a dot over a frame and a long stroke falling to the left, four strokes, as on top of \zh{房}. The next lesson puts it in your hand.

\begin{cousinweb}
\textbf{\zh{保护}} (\emph{bǎohù}) is built from \zh{保} 'to keep' and \zh{护} 'to guard', both inherited from Middle Chinese, \emph{pawX} and \emph{huH} (Guangyun, in Baxter's transcription), which glosses the first 'to keep safe' and the second 'to rescue' and 'to help'.
\end{cousinweb}

\begin{grammarlens}[title={What you've built}]
One more word for getting around.
\end{grammarlens}

\subsection*{Your turn}
\begin{itemize}
\raggedright
  \item \emph{Say it:} \emph{bǎohù}
  \item \emph{Say it:} \emph{bǎohù}, once more
  \item \emph{Say it:} \emph{wǒmen yīnggāi bǎohù héshuǐ} — We should keep the river water clean.
  \item \emph{Recall:} say the Mandarin for a great river, then the Mandarin for crystal-clear, then say \emph{bǎohù} again
\end{itemize}

\begin{tcolorbox}[breakable,colback=teal!4,colframe=teal!35!black,title={Before you move on}]
What does \zh{保护} mean? (\textquotedblleft{}To protect; to keep from harm\textquotedblright{}.) Say it once more.
\end{tcolorbox}
\section[\texorpdfstring{\zh{护} (hù)}{护 (hù)}]{\zh{护} — the hand piece, then the top of \zh{房}}

\label{lesson:ZH-W443-hu-guard}



\begin{quote}
The word first, then the character inside it.

\begin{itemize}
\raggedright
  \item \emph{Recall:} say \emph{bǎohù} — \textbf{R1}, one lesson back, before it decays
\end{itemize}

\begin{quote}
\zh{打} \zh{房}
\end{quote}

\emph{Read it:} each one aloud, with nothing said first — \emph{dǎ}, \emph{fáng}
\end{quote}

\subsection*{Script — the hand piece, then the top of \zh{房}}
\hlblockfigure{\detokenize{figures/ZH-W443-hu-guard-filmstrip.pdf}}{How \zh{护} is written, stroke by stroke}

\begin{quote}
\zh{护}
\end{quote}

\textbf{Two pieces, left before right.} Count the pieces, not the strokes.

\textbf{The left piece, the hand piece, three strokes}, as in \zh{打}: a bar, a vertical with a hook, and a stroke that rises to the right through the vertical.

\textbf{The right piece, four strokes}, as on top of \zh{房}: a dot; a stroke that runs right and turns down; a bar that closes the frame at the foot; then a long stroke falling to the left from the top corner of the frame.

3 + 4 = seven.

\subsection*{Your turn}
\begin{itemize}
\raggedright
  \item \emph{Write it:} \zh{护}, all seven, with the model in front of you
  \item \emph{Cover:} write it again from memory, and count the pieces, not the strokes
\end{itemize}

\begin{tcolorbox}[breakable,colback=teal!4,colframe=teal!35!black,title={Before you move on}]
How many strokes in \zh{护}? (\textbf{Seven}.) What is the left piece? (\textbf{The hand piece, as in \zh{打}}.) Where have you made the right piece before? (\textbf{On top of \zh{房}}.)

Source: \href{https://raw.githubusercontent.com/chanind/hanzi-writer-data/68d10a4b21150cae5e1ebbd223eed289cf32d90c/data/\%E6\%8A\%A4.json}{Hanzi Writer Data \zh{护}, pinned stroke paths and medians}.
\end{tcolorbox}
\section[\texorpdfstring{\zh{保护区} (bǎohùqū)}{保护区 (bǎohùqū)}]{\zh{保护区} (bǎohùqū) — nature reserve; protected area}

\label{lesson:ZH-C443-baohuqu}



\begin{quote}
Before the new one: say the Mandarin for crystal-clear, then the Mandarin for to keep from harm.
\end{quote}

\subsection*{You'll want to know: \zh{保护区}}
\textbf{\zh{保护区}} — \emph{bǎohùqū} — \textquotedblleft{}nature reserve; protected area\textquotedblright{}.

\begin{quote}
\textbf{\zh{我们坐车去山里的保护区}} \emph{wǒmen zuò chē qù shān lǐ de bǎohùqū} — We're driving to the nature reserve in the mountains.
\end{quote}

\begin{cousinweb}
\textbf{\zh{保护区}} (\emph{bǎohùqū}) is built from \zh{保护} 'to protect' and \zh{区} 'an area', all inherited from Middle Chinese, \emph{pawX}, \emph{huH} and \emph{khju} (Guangyun, in Baxter's transcription).
\end{cousinweb}

\begin{grammarlens}[title={What you've built}]
One more word for getting around.
\end{grammarlens}

\subsection*{Your turn}
\begin{itemize}
\raggedright
  \item \emph{Say it:} \emph{bǎohùqū}
  \item \emph{Say it:} \emph{bǎohùqū}, once more
  \item \emph{Say it:} \emph{wǒmen zuò chē qù shān lǐ de bǎohùqū} — We're driving to the nature reserve in the mountains.
  \item \emph{Recall:} say the Mandarin for crystal-clear, then the Mandarin for to keep from harm, then say \emph{bǎohùqū} again
\end{itemize}

\begin{tcolorbox}[breakable,colback=teal!4,colframe=teal!35!black,title={Before you move on}]
What does \zh{保护区} mean? (\textquotedblleft{}Nature reserve; protected area\textquotedblright{}.) Say it once more.
\end{tcolorbox}
\section[\texorpdfstring{\zh{护送} (hùsòng)}{护送 (hùsòng)}]{\zh{护送} (hùsòng) — to escort}

\label{lesson:ZH-C443-husong}



\begin{quote}
Before the new one: say the Mandarin for to keep from harm, then the Mandarin for a nature reserve.
\end{quote}

\subsection*{You'll want to know: \zh{护送}}
\textbf{\zh{护送}} — \emph{hùsòng} — \textquotedblleft{}to escort\textquotedblright{}.

\begin{quote}
\emph{tiān hēi le cūnzhǎng hùsòng wǒmen guò hé} — It got dark, so the village head saw us safely across the river.
\end{quote}

\begin{cousinweb}
\textbf{\zh{护送}} (\emph{hùsòng}) is built from \zh{护} 'to guard' and \zh{送} 'to see someone off', both inherited from Middle Chinese, \emph{huH} and \emph{suwngH} (Guangyun, in Baxter's transcription), which glosses the second 'to send'.
\end{cousinweb}

\begin{grammarlens}[title={What you've built}]
One more word for getting around.
\end{grammarlens}

\subsection*{Your turn}
\begin{itemize}
\raggedright
  \item \emph{Say it:} \emph{hùsòng}
  \item \emph{Say it:} \emph{hùsòng}, once more
  \item \emph{Say it:} \emph{tiān hēi le cūnzhǎng hùsòng wǒmen guò hé} — It got dark, so the village head saw us safely across the river.
  \item \emph{Recall:} say the Mandarin for to keep from harm, then the Mandarin for a nature reserve, then say \emph{hùsòng} again
\end{itemize}

\begin{tcolorbox}[breakable,colback=teal!4,colframe=teal!35!black,title={Before you move on}]
What does \zh{护送} mean? (\textquotedblleft{}To escort\textquotedblright{}.) Say it once more.
\end{tcolorbox}
\section[bǔyú]{bǔyú — to catch fish (for a living)}

\label{lesson:ZH-C443-buyu}



\begin{quote}
Before the new one: say the Mandarin for a nature reserve, then the Mandarin for to escort someone.
\end{quote}

\subsection*{You'll want to know: bǔyú}
\emph{bǔyú} — \textquotedblleft{}to catch fish (for a living)\textquotedblright{}.

\begin{quote}
\emph{cūnmín měitiān zuò chuán qù bǔyú} — The villagers go out fishing by boat every day.
\end{quote}

Its characters come later; for now it is a word for the ear.

\begin{cousinweb}
\emph{bǔyú} is built from \emph{bǔ} 'to seize' and \emph{yú} 'a fish', both inherited from Middle Chinese, \emph{buH} and \emph{ngjo} (Guangyun, in Baxter's transcription), which glosses the first 'to seize'; the falling-rising bǔ is irregular beside it.
\end{cousinweb}

\begin{grammarlens}[title={What you've built}]
One more word for getting around.
\end{grammarlens}

\subsection*{Your turn}
\begin{itemize}
\raggedright
  \item \emph{Say it:} \emph{bǔyú}
  \item \emph{Say it:} \emph{bǔyú}, once more
  \item \emph{Say it:} \emph{cūnmín měitiān zuò chuán qù bǔyú} — The villagers go out fishing by boat every day.
  \item \emph{Recall:} say the Mandarin for a nature reserve, then the Mandarin for to escort someone, then say \emph{bǔyú} again
\end{itemize}

\begin{tcolorbox}[breakable,colback=teal!4,colframe=teal!35!black,title={Before you move on}]
What does bǔyú mean? (\textquotedblleft{}To catch fish (for a living)\textquotedblright{}.) Say it once more.
\end{tcolorbox}
